Some efforts have sought to address multiple stages of the open-world agent loop in an integrated manner. World models \citep{ha2018world, hafner2023mastering} offer a potential unifying framework by learning environment dynamics that can serve both perception and planning, though they typically do not handle discrete, qualitative changes in the environment. End-to-end foundation models for robotics \citep{pmlr-v229-zitkovich23a, team2024octo} implicitly combine perception, reasoning, and action selection in a single model, but their behavior under genuine novelty is not well understood. Task and motion planning systems \citep{garrett2021integrated, kaelbling2011hierarchical} integrate symbolic and geometric reasoning but inherit the brittle model assumptions of classical planning. Despite growing interest in open-world autonomy, the field remains largely fragmented: methods that address perception under novelty, planning under incomplete models, and recovery from execution failures are typically developed in isolation from one another.

% =============================================================
% Thesis Contributions: Three Tenets
% =============================================================

This dissertation identifies three tenets on which progress is needed to advance the development of autonomous agents and robots for open-world environments. First, the field requires standardized environments and benchmarks that allow systematic evaluation of agents under controlled novelty, enabling rigorous comparison of approaches and driving progress in open-world AI. Second, there is a need for architectures and methods that integrate planning and learning specifically for novelty handling, allowing agents to detect execution failures, invoke learning-based recovery, and reuse acquired knowledge across subsequent encounters with novelty. Third, robotic systems require skill learning approaches that produce generalizable interaction policies, enabling robots to perform diverse manipulation tasks across objects and platforms without retraining from scratch.

These three tenets correspond to the core contributions of this dissertation. Chapter~3 introduces benchmark environments for evaluating open-world agents under controlled novelty injection. Chapter~4 develops hybrid planning-learning architectures that enable agents to adapt and recover when previously learned models become insufficient. Chapter~5 presents object-centric, force-based reinforcement learning methods for robotic manipulation that generalize across objects and robot embodiments. Together, these contributions advance the goal of building autonomous agents and robots that can operate robustly in open-world environments where novelty is not the exception but the norm.